\section*{Abstract}
Artificial intelligence already improves some forms of intellectual work, but the size and meaning of the benefit depend on what is measured. Experiments demonstrate gains in bounded writing, customer-support, and programming tasks. Observational studies increasingly show that greater artifact production can coexist with review bottlenecks, modest changes in delivered output, and uncertain commercial value. In education, assisted performance and durable learning can move in different directions. In science, specialized prediction, evaluator-guided algorithm search, and selected research workflows provide credible benefits, while reliable autonomous innovation remains uneven. This living critical narrative review connects these findings across productivity, engineering, learning, scientific discovery, business, and employment. Its central proposition is that AI's contribution should be evaluated across a complete pathway: from proposed work to verified output, use, and consequential value. The accompanying evidence map distinguishes source identity verification from substantive appraisal, and reports study design, publication status, limitations, and version information. The October 2026 release covers sources available by 12 October 2026; it does not offer a pooled effect estimate or claim exhaustive systematic coverage.


